% ============================================================================
\chapter{Implementation Considerations}
\label{ch:implementation}
% ============================================================================

This chapter discusses how to implement the hybrid architecture in
Rust, covering data structures, performance optimisations, and the
key design decisions for a proof-of-concept implementation.


% ----------------------------------------------------------------------------
\section{Why Rust?}
\label{sec:why-rust}
% ----------------------------------------------------------------------------

The choice of Rust is motivated by three factors:

\begin{enumerate}
  \item \textbf{Performance:} Rust compiles to native code with
    zero-cost abstractions. Bitwise operations map directly to
    hardware instructions with no interpreter or runtime overhead.

  \item \textbf{Safety:} Rust's ownership system prevents data races,
    buffer overflows, and memory leaks at compile time. For a
    research implementation with complex bit manipulation, this
    catches bugs early.

  \item \textbf{Portability:} Rust runs on every platform: x86,
    ARM, RISC-V, WebAssembly. The same code can target a server,
    a laptop, a Raspberry Pi, or a browser.
\end{enumerate}

\begin{engineernote}
Rust's ecosystem includes excellent support for SIMD intrinsics
(\texttt{std::arch}), which are essential for the bitwise parallel
operations in clause evaluation. The \texttt{packed\_simd} or
\texttt{std::simd} APIs allow writing architecture-aware bitwise
code that exploits 256-bit or 512-bit registers.
\end{engineernote}


% ----------------------------------------------------------------------------
\section{Data Representation}
\label{sec:data-repr}
% ----------------------------------------------------------------------------

\subsection{Bit Vectors}

The fundamental data structure is the \concept{packed bit vector}:
an array of 64-bit unsigned integers where each bit represents one
Boolean value.

For 1568 literals: $\lceil 1568 / 64 \rceil = 25$ \texttt{u64} words.

Key operations:
\begin{itemize}
  \item \textbf{AND:} Bitwise AND of two packed vectors (25
    instructions).
  \item \textbf{Equality check:} Compare two packed vectors word by
    word (25 comparisons with early exit).
  \item \textbf{Popcount:} Sum \texttt{count\_ones()} across all
    words.
  \item \textbf{NOT:} Bitwise NOT of each word (25 instructions).
\end{itemize}

\subsection{Clause Representation}

Each clause is stored as a single packed bit vector (the include
mask). Clause evaluation:

\begin{verbatim}
fn clause_matches(mask: &[u64; 25], input: &[u64; 25]) -> bool {
    for i in 0..25 {
        if (mask[i] & input[i]) != mask[i] {
            return false;
        }
    }
    true
}
\end{verbatim}

The early exit on the first non-matching word makes this very fast
for clauses that do not match (which is the common case due to
sparsity).

\subsection{Automaton States}

During training, each automaton stores its state as a \texttt{u8}
(for $N \leq 128$, states 1--256). The automata for one clause are
stored in a contiguous array of 1568 \texttt{u8} values.

\subsection{Gate Network}

Each gate stores:
\begin{itemize}
  \item Gate type: 4 bits (one of 16 types), packed as \texttt{u8}.
  \item Input indices: two \texttt{u16} values pointing to the
    previous layer.
\end{itemize}

Total per gate: 5 bytes. For 16,000 gates: 80 kB.


% ----------------------------------------------------------------------------
\section{Performance-Critical Paths}
\label{sec:critical-paths}
% ----------------------------------------------------------------------------

\subsection{Clause Evaluation (Inference Bottleneck)}

Clause evaluation dominates inference time. Optimisations:

\begin{enumerate}
  \item \textbf{SIMD:} Use 256-bit AVX2 or 512-bit AVX-512
    instructions to process 4 or 8 words per instruction instead
    of 1. This reduces the inner loop from 25 iterations to 4--7.

  \item \textbf{Early exit:} Most clauses do not match any given
    input. Checking the first few words (which often contain the
    most discriminative literals) allows early termination for
    non-matching clauses.

  \item \textbf{Cache layout:} Store clause masks contiguously in
    memory so that iterating over clauses has good spatial locality.
\end{enumerate}

\subsection{Automaton Updates (Training Bottleneck)}

Training is dominated by automaton state updates. Each update
requires:
\begin{enumerate}
  \item A random number (for stochastic feedback).
  \item A state increment or decrement.
  \item Clamping to the valid range.
\end{enumerate}

Optimisations:
\begin{itemize}
  \item Use a fast PRNG (e.g., xoshiro256 or PCG) that produces
    64 random bits per call, providing 64 random decisions at once.
  \item Vectorise the state updates using SIMD on the \texttt{u8}
    state arrays.
\end{itemize}


% ----------------------------------------------------------------------------
\section{Module Structure}
\label{sec:modules}
% ----------------------------------------------------------------------------

A suggested Rust project structure:

\begin{verbatim}
bitwise-digits/
  src/
    lib.rs          -- Public API
    bitvec.rs       -- Packed bit vector operations
    booleanise.rs   -- Thresholding + literal expansion
    tsetlin/
      mod.rs        -- TM public interface
      automaton.rs  -- Single automaton state machine
      clause.rs     -- Clause evaluation and storage
      feedback.rs   -- Type I and Type II feedback logic
      machine.rs    -- Full TM: clauses + voting
    lgn/
      mod.rs        -- LGN public interface
      gate.rs       -- Gate types and truth tables
      network.rs    -- Layered gate network
      train.rs      -- Softmax relaxation + annealing
    lut/
      mod.rs        -- LUT public interface
      spline.rs     -- B-spline evaluation
      compile.rs    -- Spline-to-LUT compilation
      score.rs      -- Scoring with compiled LUTs
    pipeline.rs     -- End-to-end inference
    mnist.rs        -- MNIST data loading
  benches/
    inference.rs    -- Inference benchmarks
    training.rs     -- Training benchmarks
  tests/
    integration.rs  -- End-to-end accuracy tests
\end{verbatim}


% ----------------------------------------------------------------------------
\section{Training Pipeline}
\label{sec:training-pipeline}
% ----------------------------------------------------------------------------

\begin{algorithm}[htbp]
\caption{Complete training pipeline.}
\label{alg:full-training}
\begin{algorithmic}[1]
  \State \textbf{// Phase 1: Train Tsetlin Machine}
  \State Load MNIST training images
  \State Booleanise all images (threshold + expand)
  \State Initialise TM~\cite{granmo2018tsetlin}: 500 clauses/class, $s = 5.0$, $T = 50$
  \For{epoch $= 1$ to $50$}
    \For{each image $(\bvec{x}, y)$ in random order}
      \State Evaluate clauses, compute votes
      \State Apply Type~I/II feedback
    \EndFor
    \State Evaluate test accuracy
  \EndFor
  \State Freeze TM clause masks
  \Statex
  \State \textbf{// Phase 2: Train Logic Gate Network}
  \State Precompute clause outputs for all training images
  \State Initialise LGN~\cite{petersen2022difflogic}: 4 layers, 4000 gates/layer, random wiring
  \For{epoch $= 1$ to $100$}
    \For{each mini-batch of clause output vectors}
      \State Forward pass with softmax relaxation
      \State Compute cross-entropy loss
      \State Backward pass, update logits
    \EndFor
    \State Anneal temperature
  \EndFor
  \State Discretise: fix each gate to its argmax type
  \Statex
  \State \textbf{// Phase 3: Compile LUTs}
  \State Compute popcount vectors for training set
  \State Train B-spline scoring functions (optional)
  \State Evaluate splines at all integer inputs
  \State Quantise to int8/int16 lookup tables
  \State Save complete model
\end{algorithmic}
\end{algorithm}


% ----------------------------------------------------------------------------
\section{Testing and Validation}
\label{sec:testing}
% ----------------------------------------------------------------------------

\paragraph{Unit tests.}
\begin{itemize}
  \item Bit vector operations: AND, OR, NOT, popcount on known inputs.
  \item Automaton transitions: verify reward/penalty move states
    correctly.
  \item Clause evaluation: manually constructed clauses on known
    images.
  \item Gate evaluation: all 16 gate types on all 4 input
    combinations.
\end{itemize}

\paragraph{Integration tests.}
\begin{itemize}
  \item Train a TM on a tiny subset (100 images) and verify it
    learns non-trivial clauses.
  \item Train an LGN on synthetic clause outputs and verify gate
    type convergence.
  \item End-to-end: train on full MNIST and verify test accuracy
    exceeds 95\%.
\end{itemize}

\paragraph{Benchmarks.}
\begin{itemize}
  \item Single-image inference latency.
  \item Full test-set evaluation throughput (images/second).
  \item Training time per epoch for each phase.
  \item Comparison with a reference FP implementation (e.g., a
    simple PyTorch 2-layer network).
\end{itemize}


% ----------------------------------------------------------------------------
\section{Benchmarking Against the Floating-Point Baseline}
\label{sec:benchmarking}
% ----------------------------------------------------------------------------

The FP baseline for comparison:

\begin{itemize}
  \item \textbf{Architecture:} 2-layer MLP, 784-300-10, ReLU
    activation.
  \item \textbf{Training:} SGD with learning rate 0.01, 50 epochs.
  \item \textbf{Expected accuracy:} $\sim$98\% on MNIST~\cite{lecun1998mnist}.
  \item \textbf{Implementation:} Either in Rust (using
    \texttt{ndarray}) or Python/PyTorch for reference.
\end{itemize}

The benchmark should measure:
\begin{enumerate}
  \item \textbf{Accuracy:} Test set accuracy for both models.
  \item \textbf{Inference speed:} Time per image, averaged over the
    10,000 test images.
  \item \textbf{Model size:} Total bytes of the saved model.
  \item \textbf{Training time:} Wall-clock time for complete training.
\end{enumerate}

\begin{engineernote}
For a fair speed comparison, both implementations should be optimised:
the FP baseline should use BLAS-accelerated matrix operations, and
the hybrid should use SIMD bitwise operations. Comparing unoptimised
implementations would be misleading in either direction.
\end{engineernote}


% ----------------------------------------------------------------------------
\section{Summary}
\label{sec:ch14-summary}
% ----------------------------------------------------------------------------

The implementation of the hybrid architecture in Rust is
straightforward because the core operations (bitwise AND, integer
increment, table lookup) map directly to hardware instructions.
The main engineering challenges are:

\begin{enumerate}
  \item Efficient packed bit vector operations with SIMD.
  \item Fast random number generation for stochastic feedback.
  \item Clean separation of training phases (integer TM, FP LGN,
    one-time LUT compilation).
  \item Thorough testing at every level to catch bit-manipulation bugs.
\end{enumerate}

With these components in place, the proof of concept can validate the
hybrid architecture's accuracy, speed, and energy promises on real
hardware.
